
\section{Additional Rules}

Last week I mentioned that we will add some of the valid arguments we prove
to our rule system. We start by defining one that we proved last week,
\emph{disjunctive syllogism}.

\begin{fitchrule}[Disjunctive Syllogism]
  \label{def:disjunctive-syllogism-rule}
  \leavevmode
  % \vspace{-0.9\baselineskip}
  There are two rules for disjunctive syllogism, one for each disjunct:\\
  \twoitems{  
    \begin{fitchproof}
      \have[m] {1} {\ephi\eor\epsi}
      \have[n] {2} {\enot\ephi}
      \have[ ] {3} {\epsi}          \by{DS}{1,2}
    \end{fitchproof}
  }{
    \begin{fitchproof}
      \have[m] {1} {\ephi\eor\epsi}
      \have[n] {2} {\enot\epsi}
      \have[ ] {3} {\ephi}          \by{DS}{1,2}
    \end{fitchproof}
  }
\end{fitchrule}

\subsection{Rules for Negation}

Recall our oft-repeated tautologies of negation.

\begin{center}
    \textbf{LEM:} $\ephi\eor\enot\ephi$         \\
    \textbf{LNC:} $\enot(\ephi\eand\enot\ephi)$ \\
    \textbf{DN:} $\enot\enot\ephi\eiff\ephi$    \\
\end{center}

We can now prove that they are all theorems.

\begin{center}
  \eproves\textbf{LEM}
  \eproves\textbf{LNC}
  \eproves\textbf{DNE}
\end{center}


\begin{table}
  %\centering
  \label{proof:lem}
  \begin{fitchproof}
    \open
    \hypo {1} {\enot(\ephi\eor\enot\ephi)} \by{AS}{}
    \open
    \hypo {2} {\ephi}                      \by{AS}{}
    \have {3} {\ephi\eor\enot\ephi}        \oi{2}
    \have {4} {\ebot}                      \ne{1,3}
    \close
    \have {5} {\enot\ephi}                 \ni{2-4}
    \have {6} {\ephi\eor\enot\ephi}        \oi{5}
    \have {7} {\ebot}                      \ne{1,6}
    \close
    \have {8} {\ephi\eor\enot\ephi}        \by{IP}{1-}
  \end{fitchproof}
  \caption{Proof of \textbf{LEM}}
\end{table}

\begin{table}
  %\centering
  \label{proof:dne}
  \begin{fitchproof}
    \open
    \hypo {1} {\ephi\eand\enot\ephi}        \by{AS}{}
    \have {2} {\ephi}                       \ae{1}
    \have {3} {\enot\ephi}                  \ae{1}
    \have {4} {\ebot}                       \ne{2,3}
    \close
    \have {5} {\enot(\ephi\eand\enot\ephi)} \ni{1-4}
  \end{fitchproof}
  \caption{Proof of \textbf{LNC}}
\end{table}

\begin{table}
  %\centering
  \label{proof:dne}
  \begin{fitchproof}
    \open
    \hypo {1} {\enot\enot\ephi} \by{AS}{}
    \open
    \hypo {2} {\enot\ephi}      \by{AS}{}
    \have {3} {\ebot}           \ne{1,2}
    \close
    \have {4} {\ephi}           \by{IP}{2-3}
    \close
    \open
    \hypo {5} {\ephi}           \by{AS}{}
    \open
    \hypo {6} {\enot\ephi}      \by{AS}{}
    \have {7} {\ebot}           \ne{1,6}
    \close
    \have {8} {\enot\enot\ephi} \ni{6-7}
    \close
    \have {9} {\ephi\eiff\enot\enot\ephi} \ei{1-4,5-8}
  \end{fitchproof}
  \caption{Proof of \textbf{DNE}}
\end{table}

Note that non-contradiction is already implemented as the negation rules. We can
add the other two theorems as derived rules.

\begin{fitchrule}[Excluded Middle]
  \label{def:excluded-middle-rule}
  \leavevmode
  \vspace{-0.9\baselineskip}
  \begin{fitchproof}
    \open
    \hypo[i] {1} {\ephi}      \by{AS}{}
    \have[j] {2} {\epsi}
    \close
    \open
    \hypo[k] {3} {\enot\ephi} \by{AS}{}
    \have[l] {4} {\epsi}
    \close
    \have[ ] {5} {\epsi}      \by{LEM}{1-2,3-4}
  \end{fitchproof}    
\end{fitchrule}

\begin{fitchrule}[Double Negative Elimination]
  \label{def:double-negative-rule}
  \leavevmode
  \vspace{-0.9\baselineskip}
  \begin{fitchproof}
    \have[m] {1} {\enot\enot\ephi}
    \have[ ] {2} {\ephi} \by{DNE}{1}
  \end{fitchproof}
\end{fitchrule}

We define the naturual deduction rules for negation in terms of \textbf{LNC},
but the rule system for TFL could be defined using any of these three
formulas, treating the other two as derived rules.\footnote{For details about
a \textbf{DN}-based system, see \cite[][]{Smi20}.}

\subsection{Modus Tollens}

Another useful technique, called \emph{modus tollens} lets us conclude that
the antecedant of a true conditional is false when its consequent is. The
rule works much like disjunctive syllogism.

\begin{fitchrule}[Modus Tollens]
  \label{def:modus-tollens-rule}
  \leavevmode
  \vspace{-0.9\baselineskip}
  \begin{fitchproof}
    \have[m] {1} {\ephi\eif\epsi}
    \have[n] {2} {\enot\epsi}
    \have[ ] {3} {\enot\ephi}     \by{MT}{1}
  \end{fitchproof}
\end{fitchrule}

This is another proof that leverages \ebot.

\begin{fitchproof}
  \hypo {1} {\ephi\eif\epsi}  \by{PR}{}
  \hypo {2} {\enot\epsi}      \by{PR}{}
  \open
  \hypo {3} {\ephi}           \by{AS}{}
  \have {4} {\epsi}           \ie{1,3}
  \have {5} {\ebot}           \ne{2,4}
  \close
  \have {6} {\enot\ephi}      \ni{3-5}
\end{fitchproof}

\subsection{De Morgans Rules}

Finally, the \emph{De Morgan rules} tell us how to properly ditribute negation
into conjunctions and disjunctions. To justify our rules, we want to prove the
following results.

\begin{center}
  \textbf{\enot\eand}
  $\enot(\ephi\eand\epsi)\efollows\eproves(\enot\ephi\eor\enot\epsi)$ \\
  \textbf{\enot\eor}
  $\enot(\ephi\eor\epsi)\efollows\eproves(\enot\ephi\eand\enot\epsi)$
\end{center}

The first one is by far the hardest.

\begin{table}[h]
  \centering
  \begin{fitchproof}
    \hypo {1} {\enot(\ephi\eand\epsi)}           \by{PR}{}
    \open
    \hypo {2} {\enot(\enot\ephi\eor\enot\epsi)}  \by{AS}{}
    \open
    \hypo {3} {\enot\ephi}                       \by{AS}{}
    \have {4} {\enot\ephi\eor\enot\psi}          \oi{3}
    \have {5} {\ebot}                            \ne{2,4}
    \close
    \have {6} {\ephi}                            \by{IP}{3-5}
    \open
    \hypo {7} {\enot\epsi}                       \by{AS}{}
    \have {8} {\enot\ephi\eor\enot\psi}          \oi{7}
    \have {9} {\ebot}                            \ne{2,8}
    \close
    \have {10}{\epsi}                            \by{IP}{7-9}
    \have {11}{\ephi\eand\epsi}                  \ai{6,10}
    \have {12}{\ebot}                            \ne{1,11}
    \close
    \have {13}{\enot\ephi\eor\enot\epsi}         \by{IP}{2-12}
  \end{fitchproof}
  \caption{Proof of $\enot(\ephi\eand\epsi)\eproves(\enot\ephi\eor\enot\epsi)$}
\end{table}

The rest are easily solved by following the proof strategies developed over the
last few weeks.

\newpage

This one should be practically trivial by now.

\begin{table}[h]
  \centering
  \begin{fitchproof}
    \hypo {1} {\enot(\ephi\eor\epsi)}     \by{PR}{}
    \open
    \hypo {2} {\ephi}                     \by{AS}{}
    \have {3} {\ephi\eor\epsi}            \oi{2}
    \have {4} {\ebot}                     \ne{1,3}
    \close
    \have {5} {\enot\ephi}                \ni{2-4}
    \open
    \hypo {6} {\epsi}                     \by{AS}{}
    \have {7} {\ephi\eor\epsi}            \oi{6}
    \have {8} {\ebot}                     \ne{1,7}
    \close
    \have {9} {\enot\epsi}                \ni{6-8}    
    \have {10}{\enot\ephi\eand\enot\epsi} \ai{5,9}
  \end{fitchproof}  
  \caption{Proof of $\enot(\ephi\eor\epsi)\eproves(\enot\ephi\eand\enot\epsi)$}
\end{table}

\newpage

These two also use the rule for \emph{disjunctive syllogism}.

\begin{table}[h]
  \centering
  \begin{fitchproof}
    \hypo {1} {\enot\ephi\eor\enot\epsi} \by{PR}{}
    \open
    \hypo {2} {\ephi\eand\epsi}          \by{AS}{}
    \open
    \hypo {3} {\enot\ephi}               \by{AS}{}
    \have {4} {\ephi}                    \ae{2}
    \have {5} {\ebot}                    \ne{3,4}
    \close    
    \open
    \hypo {6} {\enot\epsi}               \by{AS}{}
    \have {7} {\epsi}                    \ae{2}    
    \have {8} {\ebot}                    \ne{6,7}
    \close
    \have {9} {\ebot}                    \by{DS}{1,3-5,6-8}    
    \close
    \have {10}{\enot(\ephi\eand\ephi)}   \ni{2-9}
  \end{fitchproof}  
  \caption{Proof of $(\enot\ephi\eor\enot\epsi)\eproves\enot(\ephi\eand\ephi)$}
\end{table}

\newpage

\begin{table}[h]
  \centering
  \begin{fitchproof}
    \hypo {1} {\enot\ephi\eand\enot\epsi} \by{PR}{}
    \open
    \hypo {2} {\ephi\eor\epsi}            \by{AS}{}
    \open
    \hypo {3} {\ephi}                     \by{AS}{}
    \have {4} {\enot\ephi}                \ae{1}
    \have {5} {\ebot}                     \ne{3,4}
    \close    
    \open
    \hypo {6} {\epsi}                     \by{AS}{}
    \have {7} {\enot\epsi}                \ae{1}
    \have {8} {\ebot}                     \ne{6,7}
    \close
    \have {9} {\ebot}                     \by{DS}{2,3-5,6-8}
    \close
    \have {10}{\enot(\ephi\eor\epsi)}     \ni{2-9}
  \end{fitchproof}  
  \caption{Proof of $(\enot\ephi\eand\enot\epsi)\eproves\enot(\ephi\eor\epsi)$}
\end{table}

Once again, we can incorporate these results with a set of explicit rules.

\Needspace{12\baselineskip}
\begin{fitchrule}[De Morgan Rules]
  \label{def:de-morgan-rules}
  \leavevmode
  % \vspace{-0.9\baselineskip}
  There are four rules for handling negated coordinations forming a square of
  duals and rule reversals.

  We have a pair of rules for negated conjunctions (\enot\eand): \\
  \twoitems{  
    \begin{fitchproof}
      \have[m] {1} {\enot(\ephi\eand\epsi)}
      \have[ ] {2} {\enot\ephi\eor\enot\epsi}   \by{DeM}{1}
    \end{fitchproof}
  }{
    \begin{fitchproof}
      \have[m] {1} {\enot\ephi\eor\enot\epsi}
      \have[ ] {2} {\enot(\ephi\eand\epsi)}     \by{DeM}{1}
    \end{fitchproof}
  }

  And a dual pair of rules for negated disjunctions (\enot\eor): \\
  \twoitems{  
    \begin{fitchproof}
      \have[m] {1} {\enot(\ephi\eor\epsi)}
      \have[ ] {2} {\enot\ephi\eand\enot\epsi}  \by{DeM}{1}
    \end{fitchproof}
  }{
    \begin{fitchproof}
      \have[m] {1} {\enot\ephi\eand\enot\epsi}
      \have[ ] {2} {\enot(\ephi\eor\epsi)}      \by{DeM}{1}
    \end{fitchproof}
  }
\end{fitchrule}

